% GENERATED FILE. Edit canonical lessons, then run npm run generate:books.
% canonical-source-hash: fnv1a64:dd1862a9e93ba3a2
% canonical-lessons: RU-C269-yesli, RU-C269-yesli-budet, RU-C269-zavisit, RU-C269-review

\chapter{If, and It Depends}
\label{ch:ru-if-and-it-depends}
\hlchaptermodality{269}

\begin{chapteropening}
\textbf{By the end of this chapter:} \emph{I can say what will happen if something happens, with \ru{если} and the future in both halves, and answer a question with \ru{это} \ru{зависит} \ru{от} ...}

Complete RU-C269-review: a weekend plan with \ru{если} and the future, a whether-question with \ru{ли}, and an answer that depends.
\end{chapteropening}

\section[\texorpdfstring{\ru{если} …, … (yésli …, …)}{если …, … (yésli …, …)}]{\ru{если} — if, and the if that is a question}

\label{lesson:RU-C269-yesli}



\begin{quote}
Say \textquotedblleft{}if\textquotedblright{}. (\emph{Yésli.}) Then say \textquotedblleft{}I've been learning Russian for two years\textquotedblright{}. (\emph{Ya uzhé dva góda uchú rússkiy.}) The word for \textquotedblleft{}if\textquotedblright{} has waited since the joining words of the A2 chapters. Now it joins.
\end{quote}

\begin{grammarlens}[title={Grammar Lens: a condition and a comma}]
\textbf{\ru{Если}} works like \emph{\ru{когда}}: it hangs one half on the other, and the comma sits at the seam, whichever half comes first.

\begin{quote}
\textbf{\ru{Если} \ru{хочешь}, \ru{мы} \ru{можем} \ru{пойти} \ru{в} \ru{кино}.} — \emph{yésli khóchesh, my mózhem poytí v kinó} — If you like, we can go to the cinema. \textbf{\ru{Мы} \ru{можем} \ru{пойти} \ru{в} \ru{кино}, \ru{если} \ru{хочешь}.} — the same, turned round. \textbf{\ru{Если} \ru{ты} \ru{устал}, \ru{отдохни}.} — \emph{yésli ty ustál, otdokhní} — If you're tired, have a rest.
\end{quote}

\emph{\ru{Пойти}} is \textquotedblleft{}to go, to set off\textquotedblright{}, the partner of \emph{\ru{идти}}; \emph{\ru{отдохни}} is \emph{\ru{отдохнуть}} to a friend.
\end{grammarlens}

\begin{grammarlens}[title={Grammar Lens: \textquotedblleft{}if\textquotedblright{} that means \textquotedblleft{}whether\textquotedblright{}}]
English \textquotedblleft{}if\textquotedblright{} has a second job: \textquotedblleft{}I don't know if the museum is open\textquotedblright{}. That \textquotedblleft{}if\textquotedblright{} is a question, \textquotedblleft{}whether\textquotedblright{}, and Russian answers it with \textbf{\ru{ли}}, not \emph{\ru{если}}:

\begin{quote}
\textbf{\ru{Я} \ru{не} \ru{знаю}, \ru{работает} \ru{ли} \ru{музей}.} — \emph{ya ne znáyu, rabótayet li muzéy} — I don't know if the museum is open.
\end{quote}

The verb comes first, then \textbf{\ru{ли}}: \textquotedblleft{}is-it-working, the museum\textquotedblright{}. A museum that is open, in Russian, \textquotedblleft{}works\textquotedblright{}.
\end{grammarlens}

\begin{cousinweb}
\textbf{\ru{если}} was once \emph{\ru{есть} \ru{ли}}, \textquotedblleft{}is it?\textquotedblright{}: \emph{\ru{есть}}, \textquotedblleft{}is\textquotedblright{}, with the asking word \emph{\ru{ли}}. EtymDB 2.1 records it as a compound of \emph{\ru{есть}} and \emph{\ru{ли}}, though it glosses that \emph{\ru{есть}} as \textquotedblleft{}to eat\textquotedblright{}, the other word with the same spelling; \textquotedblleft{}is\textquotedblright{}, the form of \emph{\ru{быть}}, is the one that fits. A condition began as a question.
\end{cousinweb}

\subsection*{Your turn}
Say these aloud:

\begin{itemize}
\raggedright
  \item \emph{Yésli khóchesh, my mózhem poytí v kinó.}
  \item the same, turned round: \emph{My mózhem poytí v kinó, yésli khóchesh.}
  \item \emph{Yésli ty ustál, otdokhní.}
  \item \emph{Ya ne znáyu, rabótayet li muzéy.}
\end{itemize}

\begin{tcolorbox}[breakable,colback=teal!4,colframe=teal!35!black,title={Before you move on}]
Where does the comma go with \emph{\ru{если}}? (\textbf{At the seam}, either order.) \textquotedblleft{}I don't know if he understands\textquotedblright{}: \emph{\ru{если}} or \emph{\ru{ли}}? (\emph{\textbf{\ru{Ли}:}} \emph{\textbf{\ru{Я} \ru{не} \ru{знаю}, \ru{понимает} \ru{ли} \ru{он}.}})
\end{tcolorbox}
\section[\texorpdfstring{\ru{если} \ru{будет} … (yésli búdet …)}{если будет … (yésli búdet …)}]{\ru{если} \ru{будет} \ru{дождь} — the future on both sides of if}

\label{lesson:RU-C269-yesli-budet}



\begin{quote}
Say \textquotedblleft{}if you like, we can go to the cinema\textquotedblright{}. (\emph{Yésli khóchesh, my mózhem poytí v kinó.}) Then \textquotedblleft{}I hope the weather will be good\textquotedblright{}. (\emph{Ya nadéyus, shto pogóda búdet khoróshaya.}) That \emph{\ru{будет}} is the key to this lesson.
\end{quote}

\begin{grammarlens}[title={Grammar Lens: the future on both sides}]
\begin{quote}
\textbf{\ru{Если} \ru{в} \ru{субботу} \ru{будет} \ru{дождь}, \ru{мы} \ru{пойдём} \ru{в} \ru{музей}.} — \emph{yésli v subbótu búdet dozhd, my poydyóm v muzéy} — If it rains on Saturday, we'll go to the museum.
\end{quote}

English puts the \textquotedblleft{}if\textquotedblright{} half in the present: \textquotedblleft{}if it rains\textquotedblright{}. When the condition lies ahead, Russian usually puts it in the future as well: \textquotedblleft{}if there will be rain\textquotedblright{}.

\begin{quote}
\textbf{\ru{Если} \ru{будет} \ru{время}, \ru{я} \ru{позвоню}.} — \emph{yésli búdet vrémya, ya pozvonyú} — If I have time, I'll ring.
\end{quote}
\end{grammarlens}

\begin{grammarlens}[title={Grammar Lens: a perfective points forward}]
\textbf{\ru{Позвоню}}, \textbf{\ru{пойдём}}: these look like present forms, and they are future. A perfective verb has no present, because a deed seen whole cannot be under way now; its present-looking forms point to one deed to come.

\begin{quote}
\textbf{\ru{Если} \ru{я} \ru{прочитаю} \ru{книгу}, \ru{я} \ru{расскажу} \ru{тебе}.} — \emph{yésli ya prochitáyu knígu, ya rasskazhú tebé} — If I finish the book, I'll tell you.
\end{quote}

The imperfective builds its future with \emph{\ru{буду}} and the infinitive, as met beside \emph{\ru{будущий}}: \textbf{\ru{Я} \ru{буду} \ru{читать}} — I will be reading.
\end{grammarlens}

\begin{cousinweb}
\textbf{\ru{будет}} is the future of \textbf{\ru{быть}}, which EtymDB 2.1 takes back to Proto-Indo-European \emph{*b\textsuperscript{h}uH-}, \textquotedblleft{}to become\textquotedblright{}. The future of \textquotedblleft{}to be\textquotedblright{} keeps a root that meant becoming: what will be is what becomes.
\end{cousinweb}

\subsection*{Your turn}
Say these aloud:

\begin{itemize}
\raggedright
  \item \emph{Yésli v subbótu búdet dozhd, my poydyóm v muzéy.}
  \item \emph{Yésli búdet vrémya, ya pozvonyú.}
  \item \emph{Yésli ya prochitáyu knígu, ya rasskazhú tebé.}
  \item a plan for the weekend of your own, with \emph{yésli búdet} and one perfective
\end{itemize}

\begin{tcolorbox}[breakable,colback=teal!4,colframe=teal!35!black,title={Before you move on}]
\textquotedblleft{}If it rains\textquotedblright{}: present or future in Russian? (\textbf{Future:} \emph{\textbf{\ru{если} \ru{будет} \ru{дождь}}}.) Is \emph{\ru{позвоню}} now or later? (\textbf{Later}: a perfective, one call to come.)
\end{tcolorbox}
\section[\texorpdfstring{\ru{это} \ru{зависит} \ru{от} … (éto zavísit ot …)}{это зависит от … (éto zavísit ot …)}]{\ru{это} \ru{зависит} \ru{от} … (éto zavísit ot …) — it depends}

\label{lesson:RU-C269-zavisit}



\begin{quote}
Say \textquotedblleft{}if it rains on Saturday, we'll go to the museum\textquotedblright{}. (\emph{Yésli v subbótu búdet dozhd, my poydyóm v muzéy.}) Then \textquotedblleft{}the weather\textquotedblright{}. (\emph{Pogóda.}) Sometimes the honest answer to a question is a condition.
\end{quote}

\subsection*{You'll want to know: \ru{это} \ru{зависит} \ru{от} …}
\begin{quote}
\textbf{— \ru{Ты} \ru{пойдёшь} \ru{в} \ru{парк}? — \ru{Не} \ru{знаю}. \ru{Это} \ru{зависит} \ru{от} \ru{погоды}.} — \emph{ty poydyósh v park? ne znáyu. éto zavísit ot pogódy} — Are you going to the park? I don't know. It depends on the weather.
\end{quote}

\textbf{\ru{Зависит}} is \textquotedblleft{}it depends\textquotedblright{}, from \emph{\ru{зависеть}}. In quick conversation it can shrink to \textbf{\ru{Зависит}.} — \textquotedblleft{}It depends.\textquotedblright{}

\begin{grammarlens}[title={Grammar Lens: \ru{от} and the \textquotedblleft{}of\textquotedblright{} form}]
\textbf{\ru{от}} — \emph{ot} — \textquotedblleft{}from, on\textquotedblright{}. What a thing depends on takes the same form as after \emph{\ru{до}}: \emph{\ru{погода}}, \textbf{\ru{от} \ru{погоды}}; \emph{\ru{цена}}, \textbf{\ru{от} \ru{цены}}.

\begin{quote}
\textbf{\ru{Это} \ru{зависит} \ru{от} \ru{цены}.} — \emph{éto zavísit ot tsený} — It depends on the price. \textbf{\ru{Это} \ru{зависит} \ru{от} \ru{времени}.} — \emph{éto zavísit ot vrémeni} — It depends on the time. \emph{\ru{Время}} has its own shape here: \textbf{\ru{времени}}. \textbf{\ru{Это} \ru{зависит} \ru{от} \ru{тебя}.} — \emph{éto zavísit ot tebyá} — It's up to you.
\end{quote}
\end{grammarlens}

\begin{cousinweb}
\textbf{\ru{зависеть}} is \emph{\ru{за}-} with \textbf{\ru{висеть}}, \textquotedblleft{}to hang\textquotedblright{}, which EtymDB 2.1 traces through Old East Slavic \emph{visěti} to Proto-Slavic \emph{*visěti}. To depend is to hang on something. English draws the same picture: \emph{depend} is Latin \emph{dēpendēre}, \textquotedblleft{}to hang down from\textquotedblright{} (Lewis \& Short).
\end{cousinweb}

\subsection*{Your turn}
\begin{itemize}
\raggedright
  \item \emph{Say it:} \emph{Éto zavísit ot pogódy.}
  \item \emph{Say it:} \emph{Éto zavísit ot tsený}, then \emph{ot vrémeni}
  \item \emph{Answer:} \emph{Ty poydyósh v park?} \emph{Ne znáyu. Éto zavísit ot …}
  \item \emph{Say it:} \emph{Mne kázhetsya, shto éto zavísit ot tebyá.}
\end{itemize}

\begin{tcolorbox}[breakable,colback=teal!4,colframe=teal!35!black,title={Before you move on}]
\textquotedblleft{}It depends on the weather\textquotedblright{}? (\emph{\textbf{\ru{Это} \ru{зависит} \ru{от} \ru{погоды}.}}) \textquotedblleft{}It's up to you\textquotedblright{}? (\emph{\textbf{\ru{Это} \ru{зависит} \ru{от} \ru{тебя}.}}) What verb is inside \emph{\ru{зависеть}}? (\textbf{One meaning \textquotedblleft{}to hang\textquotedblright{}.})
\end{tcolorbox}
\section[(review)]{Review — if, and it depends}

\label{lesson:RU-C269-review}



\begin{quote}
Say \textquotedblleft{}if you like, we can go to the cinema\textquotedblright{}. (\emph{Yésli khóchesh, my mózhem poytí v kinó.}) Then \textquotedblleft{}I don't know if the museum is open\textquotedblright{}. (\emph{Ya ne znáyu, rabótayet li muzéy.})
\end{quote}

\begin{grammarlens}[title={What you've built}]
Three lessons, and Russian can now look ahead on a condition. \textbf{\ru{Если}} with a comma at the seam. The future on both sides when the condition lies ahead, with \textbf{\ru{будет}} and a perfective's forward-pointing forms. \textbf{\ru{Ли}} when \textquotedblleft{}if\textquotedblright{} means \textquotedblleft{}whether\textquotedblright{}. And \textbf{\ru{это} \ru{зависит} \ru{от}} with the \textquotedblleft{}of\textquotedblright{} form.
\end{grammarlens}

\subsection*{The exchange}
\begin{quote}
\textbf{— \ru{Ты} \ru{пойдёшь} \ru{в} \ru{музей} \ru{в} \ru{субботу}?} \textbf{— \ru{Это} \ru{зависит} \ru{от} \ru{погоды}. \ru{Если} \ru{будет} \ru{дождь}, \ru{пойду}.}
\end{quote}

\emph{Poydú}, \textquotedblleft{}I'll go\textquotedblright{}: the perfective, one visit to come. The answer depends, and then says on what.

\subsection*{Your turn}
\begin{itemize}
\raggedright
  \item \emph{Say it:} \emph{Éto zavísit ot pogódy. Yésli búdet dozhd, poydú.}
  \item \emph{Say it:} \emph{Yésli búdet vrémya, ya pozvonyú.}
  \item \emph{Say it:} \emph{Póyezd deshévle, chem samolyót, no samolyót bystréye.}
  \item \emph{Answer:} \emph{Ty poydyósh v park?} Make it depend on something, then add one \emph{yésli}
  \item \emph{Say it:} \emph{Dva góda nazád ya byl v Moskvé, i ya uzhé dva góda uchú rússkiy.}
\end{itemize}

\begin{tcolorbox}[breakable,colback=teal!4,colframe=teal!35!black,title={Before you move on}]
\textquotedblleft{}If there's time\textquotedblright{}, for tomorrow? (\emph{\textbf{\ru{Если} \ru{будет} \ru{время}.}}) \textquotedblleft{}Whether\textquotedblright{}: \emph{\ru{если}} or \emph{\ru{ли}}? (\emph{\textbf{\ru{Ли}.}}) \textquotedblleft{}It depends on the price\textquotedblright{}? (\emph{\textbf{\ru{Это} \ru{зависит} \ru{от} \ru{цены}.}})
\end{tcolorbox}
